%==========================================================================================
\section{Waveforms: the same incident $h_+$ pulse in the two parities}\label{sec:time}
\begin{question}
What does a detector see if the same pulse arrives once carried by the even master function and once by the odd one?
\end{question}
\paragraph{Set-up.} For a single $(\ell,m)$ the angular factor in Eq.~\eqref{eq:spin} is the same for the incident and the
reflected wave. So $h_+$ and $h_\times$ at any fixed direction are that angular factor times $\Psi_{\rm in}(v)$ before the
reflection and $\Psi_{\rm out}(u)$ after it. For instance, for $m=0$ the even channel is literally ``incident $h_+$, reflected
$h_+$'', and the odd channel is the same with $h_\times$. We therefore compare the reflected $\Psi_{\rm out}(u)$ of the two
parities for the same incident pulse,
\begin{equation}
\Psi_{\rm in}(v)=e^{-v^2/2\sigma^2}\cos(\omega_0v),\qquad \sigma=8M,\quad \omega_0=0.5/M,\qquad v=t+r_* ,
\end{equation}
by Fourier synthesis, $\Psi_{\rm out}(u)=\frac1{2\pi}\int S(\omega)\tilde\Psi_{\rm in}(\omega)e^{-i\omega u}\dd\omega$, with $u=t-r_*$. The
integral is done with a Filon-type quadrature on the adaptive frequency grid, so the narrow resonances are resolved.

\paragraph{Causality and the instability.} The fluid shell has an unstable even-parity matter mode at $\omega=i\gamma$.
There, $S_{\rm even}$ has a pole in the upper half plane, and the real-axis integral is then \emph{not} the causal response.
The causal (retarded) response is obtained by moving the contour above the pole:
\begin{equation}
\Psi^{\rm causal}_{\rm out}(u)=\Psi^{\rm real\ axis}_{\rm out}(u)-i\,{\rm Res}_{\omega=i\gamma}\big[S_{\rm even}\tilde\Psi_{\rm in}\big]\,e^{\gamma u}.
\end{equation}
The residue must be computed directly. For ${\rm Im}\,\omega>0$ the outgoing solution decays with $r$, so we project the
exterior solution onto the in/out asymptotic solutions at a moderate radius. This reproduces the real-axis $S$ to
$10^{-8}$, and Newton iteration on $B_{\rm in}=0$ returns the pole exactly at the known unstable frequency:
$0.1332058i$ ($R=3M$), $0.2442921i$ ($2.2M$), $0.0401379i$ ($6M$). The residues, stable to three digits against the
projection radius, are tiny: $7.5\times10^{-5}i$, $8.2\times10^{-6}i$ and $1.1\times10^{-5}i$. The unstable mode lives in the
shell's matter and is only weakly coupled to waves coming from infinity. For this pulse its amplitude at $u=0$ is
$3.9\times10^{-7}$ ($R=3M$), $1.4\times10^{-8}$ ($2.2M$) and $2.3\times10^{-8}$ ($6M$). It reaches the amplitude of the incident
pulse at $u\approx111M$, $74M$ and $438M$ respectively.

\paragraph{Check.} For the odd channel the result was compared with a direct time-domain (leapfrog) evolution of
$\partial_t^2\Psi=\partial_x^2\Psi-V\Psi-\Delta\,\delta(x-x_R)\Psi$ that shares no code with the frequency-domain calculation. With
source and detector far from the object ($x=450M$ and $350M$), and the $O(1/\omega r)$ near-zone factors of the in/out
solutions included, the two agree to $0.4\%$ ($R=3M$) and $0.9\%$ ($2.2M$) of the peak (Fig.~\ref{fig:lf}).

\begin{figure}[p]\centering
\includegraphics[width=0.49\textwidth]{r4_f06_time_bh.pdf}\hfill
\includegraphics[width=0.49\textwidth]{r4_f07_time_R3.pdf}\\[2mm]
\includegraphics[width=0.49\textwidth]{r4_f08_time_R22.pdf}\hfill
\includegraphics[width=0.49\textwidth]{r4_f09_time_log.pdf}
\caption{Reflected waveforms of the same incident pulse, carried by the odd (blue) or even (orange) master function. Top
left: black hole. Top right: shell at $3M$. Bottom left: shell at $2.2M$. Dotted: the growing matter mode of the even
channel, added causally. Bottom right: long-time view (log scale) for $R=6M$ and $3M$. The even channel keeps ringing at
the matter frequency and is eventually overtaken by the unstable mode; the odd channel just decays. The floor at
$\sim10^{-5}$ is numerical.}\label{fig:time}
\end{figure}

\paragraph{What the waveforms show (Fig.~\ref{fig:time}).}
\begin{itemize}
\item \textbf{Black hole.} The two waveforms are almost identical. The even one is the odd one passed through the
all-pass filter $(\lambda_\ell-12M\partial_u)/(\lambda_\ell+12M\partial_u)$, the time-domain form of Eq.~\eqref{eq:bhphase}. This
gives a delay of $\simeq24M/\lambda_\ell=1M$ for $\ell=2$ at low frequency. Both ring down at the same QNM.
\item \textbf{Shell at $3M$.} The prompt reflection is again nearly the same in both parities (the retardance is only
$\sim10^\circ$ in this band). The difference is at late times: the even channel keeps a weak ringing at the matter frequency
$\omega M=0.197$, with a decay time $\sim3000M$. On top of it the unstable mode grows as $e^{0.133u/M}$ and dominates after
$\sim110M$. The odd channel has neither.
\item \textbf{Shell at $2.2M$.} Both parities ring in their trapped modes, at slightly different frequencies ($0.421$ odd,
$0.447$ even) and dampings ($0.035$ vs $0.044$). The two waveforms dephase visibly after a few cycles, and the unstable
mode takes over the even channel after $\sim74M$.
\item \textbf{Shell at $6M$.} The prompt reflections are indistinguishable. The even channel carries a long matter
ringing at $\omega M=0.089$ at the $10^{-5}$ level and blows up after $\sim440M$.
\end{itemize}

\paragraph{Polarization in the time domain.} For a pulse with equal $E$ and $B$ content, the part that comes back with the
orthogonal polarization is $(\Psi^{\rm even}_{\rm out}-\Psi^{\rm odd}_{\rm out})/2$ (Fig.~\ref{fig:rot}). It carries $3.7\%$ of the
reflected energy for the black hole and $1.0\%$ for the $R=3M$ shell (stable part), concentrated at the edges of the pulse,
where the phase difference acts. In the shell case the rotated component then persists at late times, carried by the
even-only matter ringing.

\begin{figure}[h]\centering
\includegraphics[width=0.49\textwidth]{r4_f10_converted_time.pdf}\hfill
\includegraphics[width=0.49\textwidth]{r4_f11_leapfrog_check.pdf}
\caption{Left: preserved and rotated parts of a reflected pulse with equal $E$ and $B$ content. Right: check of the
frequency-domain synthesis against an independent time-domain evolution (odd channel).}\label{fig:rot}\label{fig:lf}
\end{figure}

\begin{learned}
Plotted as $h_+(t)$, the even and odd reflections of the same pulse look alike during the prompt response. Their
difference grows with time: in the dephasing of the trapped modes, in the even-only matter ringing, and above all in the
even-only instability of the static shell.
\end{learned}
